\section{When is the first coefficient depth sufficient?}
\label{sec:first-depth}
Berger and Sprang define the numerical invariant $s_0(n)$ and compute it for a range of degrees \cite[Definition 5.1 and Section 5]{BS}. Their choice is the unramified quadratic field $K/\mathbb Q_p$, $q=p^2$, and the coordinate with $[p](X)=pX+X^q$. Write $M=\mathcal O[[X]]^{\psi=0}$ for the kernel of their normalized trace. For $f\in M$, let $c_{f,m}(T)$ be the polynomial characterized by $c_{f,m}(a)=[X^m]f([a](X))$ for $a\in\mathcal O$. Its degree is at most $m$. Then $s_0(n)$ is the first cutoff $N\ge n$ at which the $\mathcal O$-span of these polynomials with $m\le N$ contains a degree-$n$ polynomial attaining the optimal leading valuation $-w_q(n)$, where
\[
 w_q(n)=\sum_{j\ge1}\lfloor n/q^j\rfloor=(n-s_q(n))/(q-1)
\]
and $s_q(n)$ is the base-$q$ digit sum. The finiteness of $s_0(n)$ follows from the quadratic character-variety result and the criterion of \cite{BS,BSS}. The following exact statements address the size of the cutoff.

\begin{theorem}[First-depth criterion for small digit sum]\label{thm:first-depth}
Write $n=\sum_{j\ge0}n_jq^j>0$ and $s=\sum_j n_j$. If $1\le s\le q-2$, then $s_0(n)=n$ exactly when adding the integers $n_j$ in base $p$ involves no carries. If $s=q-1$, or if $s\le q-2$ and there is a carry, then $s_0(n)\ge n+q-1$. At the first excluded digit sum the specific value $s_0(q-1)=2q-2$ is exact.
\end{theorem}
\begin{proof}
Let $E(Z)=Z+\cdots$ be the Lubin--Tate exponential and $L(X)$ its compositional inverse. Then $[a](X)=E(aL(X))$. In $c_{f,n}$, the coefficient of $T^n$ is $[Z^n]f(E(Z))$, since $[X^n]L(X)^n=1$. Thus only the $n$th row can supply the optimal leading term at depth $n$.

Take $\eta^{q-1}=p$, and set $V(Z)=E(\eta Z)/\eta$. The exponential equation becomes $V(pZ)=pV(Z)+pV(Z)^q$. Its coefficients are integral: for $m>1$,
\[
 (p^{m-1}-1)[Z^m]V=[Z^m]V^q,
\]
with a unit on the left and only earlier coefficients on the right. In fact this recursion puts all coefficients in $\mathbb Z_p$, with support in degrees $1\bmod(q-1)$. Reducing yields $\overline V+\overline V^q=Z$, whence $\overline V=\sum_{j\ge0}(-1)^j Z^{q^j}$. A monomial $X^i$ of $f$ can only contribute if $i\equiv n\bmod(q-1)$, and its valuation is at least $-(n-i)/(q-1)$. When $1\le s\le q-1$, the first eligible positive $i$ is $s$. Larger $i$ have valuation strictly above $-w_q(n)$. The constant monomial contributes zero.

For $i=s$, the reduction of $\eta^{n-s}[Z^n]E^s$ is the coefficient of $Z^n$ in $\overline V^s$. Since $s<q$, writing $n$ as a sum of $s$ powers of $q$ uses $n_j$ copies of $q^j$. The coefficient is, up to sign, the multinomial $s!/\prod_j n_j!$, a unit exactly when addition is carry-free in base $p$, by Legendre's factorial formula. If $s\le q-2$, the monomial $X^s$ belongs to $M$ by \cite[Lemma 5.2]{BS}. For $s=q-1$ the same lemma gives $[X^{q-1}]f\in p\mathcal O$ for every $f\in M$: reduce its generators modulo $(p,X^q)$ to see only the constant and powers $X^1,\ldots,X^{q-2}$. This proves the criterion at depth $n$.

For any $(q-1)$st root of unity $\zeta$, $[\zeta](X)=\zeta X$ in this coordinate (it commutes with $[p]$ and has the required derivative). Consequently $c_{f,m}(\zeta T)=\zeta^m c_{f,m}(T)$. No intervening row $n<m<n+q-1$ supplies a $T^n$ term, proving the lower bound.

Finally set $m=2q-2$. The series $g=pX^{q-1}-(1-q)$ and $h=X^{q-2}[p](X)-g=X^m+(1-q)$ both belong to $M$ by \cite[Lemma 5.2]{BS}. The logarithm has coefficient $(p-p^q)^{-1}$ at $X^q$, as follows by comparing $L(pX+X^q)=pL(X)$. Therefore
\[
 c_{g,m}(T)=\frac{p(q-1)}{p-p^q}(T^{q-1}-T^m),\qquad c_{h,m}(T)=T^m.
\]
The displayed multiplier is a unit. Their span contains $T^{q-1}$, whose leading valuation is $0=-w_q(q-1)$, at the claimed cutoff.
\end{proof}
For $p=2$, $q=4$, this gives $s_0(2\cdot4^j)=2\cdot4^j$ for every $j\ge0$. It gives $s_0(4^i+4^j)\ge 4^i+4^j+3$ for $i\ne j$ since $1+1$ carries in base two. In particular $s_0(5)\ge8$ and $s_0(3)=6$. The argument isolates the first attainable leading coefficient of a formal exponential, in the same spirit as the CM-derivative extremum in the Lean source. It leaves open the exact later depth for the delayed families and the general behavior for larger digit sums. We found no matching closed criterion in the sources reviewed; the theorem makes no claim of publication priority.

\begin{corollary}[Unbounded excess depth]\label{cor:depth-gap}
Suppose $s_0(n)>n$, and let $J\ge1$ satisfy $J+\lfloor\log_p J\rfloor\le v_p(n)$. Then $s_0(n)\ge n+(q-1)(J+1)$. In particular, for every fixed prime $p$, the excess $s_0(n)-n$ is unbounded. Along $n=(q-1)q^k$ it is at least $(q-1)(2k-O(\log(k+1)))$ as $k\to\infty$.
\end{corollary}
\begin{proof}
In any row, the exact coefficient of $T^n$ in $c_{f,m}$ factors as
\[
 [Z^n]f(E(Z))\,[X^m]L(X)^n.
\]
The first factor has valuation strictly above $-w_q(n)$ for every $f\in M$ by the assumed failure at depth $n$ and the integer-valued leading-coefficient bound. We show that the second factor stays integral throughout the asserted range. The compositional inverse of the integral series $V$ with unit linear coefficient is $\eta^{-1}L(\eta X)=X(1+D(X))$ with $D\in X^{q-1}\mathcal O_{K(\eta)}[[X]]$, supported at positive multiples of $q-1$. For $m=n+j(q-1)$ the scaled coefficient $p^j[X^m]L(X)^n$ equals $[X^{j(q-1)}](1+D)^n$. If $j\ge1$, all contributions carry a factor $\binom ni$ with $1\le i\le\min(n,j)$. From $i\binom ni=n\binom{n-1}{i-1}$ their valuations are at least $v_p(n)-\lfloor\log_p j\rfloor$. Consequently $[X^m]L(X)^n$ is integral for $1\le j\le J$. The other rows have zero coefficient by the degree symmetry. On $n=(q-1)q^k$, the theorem gives the required initial failure and $v_p(n)=2k$. For all sufficiently large $k$ take $J=2k-\lfloor\log_p(2k)\rfloor$.
\end{proof}
